% 夜半用户手册 · 第 4 章：钢琴卷帘（中文）
% T-DOC-001 逐章交付（CH04 任务未到前预写）。
% 事实来源：app.slint 卷帘属性与菜单、help.rs 引导步骤、bridge.rs 音高车道映射、
% docs/manual/04-piano-roll.md、cli.rs SHORTCUTS。

\chapter{钢琴卷帘}
\label{ch:piano-roll}

\classquote{凡音之起，由人心生也。}{《礼记·乐记》}{音符落在卷帘上之前，先有想表达的情绪；技巧与规则都服务于听感。}





钢琴卷帘是写 MIDI 音符的地方：横轴是时间，纵轴是音高。这一章讲怎么打开它、
怎么画音符、怎么改力度与量化，以及当前版本在拖动上的边界。

\section{打开与视图导航}

\subsection{打开卷帘}

按 P 键，或从底部控制台第一个页签切到钢琴卷帘；F7 也是卷帘页签的直达键。
选中一条 MIDI 合成器轨上的片段，卷帘就显示那段片段里的音符。

\subsection{看清两个维度}

横轴是时间（以 tick 计，按工程拍号换算成小节与拍）；纵轴是音高，从上到下
按音高车道排列，左侧有键标签（如 C4、Eb5）作为音高参照。窗口顶部的读数会
显示当前可见的音高范围。底部是力度泳道，音符的力度用竖条表示。

\screenshot{ch04-roll-overview}{钢琴卷帘全景，标出时间标尺、音高车道与键标签、力度泳道}

\subsection{缩放与取景}

Cmd/Ctrl+↑ / ↓ 做音轨垂直缩放，Cmd/Ctrl+= / - 做时间轴水平缩放。卷帘有
「取景」概念：可以按当前片段、当前轨或全部内容来定可见范围；超出可见范围的
高音或低音不会在窗口里画出，切换取景范围即可看到。

\section{音符录入、选择与移动}

\subsection{画音符}

新手引导里说得很直白：\textbf{数字键 2 是绘制工具}，在卷帘上点一下就落一个
音符。落点由点击位置决定：横向是时间，纵向由 y 坐标反查到对应音高车道。
画完后可以拖动音符改时间与音高，拖右边缘改长度。

\subsection{选择与拖动}

点一个音符即选中它。当前版本的指针拖动只作用于指针下那一颗音符——即使你
选中了多颗，拖其中一颗，其余音符不会跟着动。因此「整组拖动」与「整组改
时值」目前不成立；要批量移动，请分次操作，或等到该能力正式提供。夜半不会
把未实现的操作伪装成成功。

\begin{figure}[htbp]
\centering
\begin{tikzpicture}[node distance=0.5cm, >=stealth,
    box/.style={draw, rounded corners, align=center, font=\small, minimum width=2.0cm, minimum height=0.9cm}]
  \node[box] (a) {选片段};
  \node[box, right=of a] (b) {数字键 2 画音符};
  \node[box, right=of b] (c) {拖动改音高/时值};
  \node[box, right=of c] (d) {量化 / 人性化};
  \node[box, right=of d] (e) {试听};
  \draw[->] (a)--(b); \draw[->] (b)--(c); \draw[->] (c)--(d); \draw[->] (d)--(e);
\end{tikzpicture}
\caption{卷帘里的常用流程：选片段 $\to$ 画音符 $\to$ 改音高/时值 $\to$ 量化 $\to$ 试听}
\label{fig:ch04-roll-flow}
\end{figure}

\section{力度、量化与音阶工具}

\subsection{力度}

底部力度泳道里，每个音符对应一根竖条，高度就是力度值。可以直接在泳道里
调整单颗音符的力度；拖动时卷帘会给出预览读数。

\subsection{量化与人性化}

「编辑 $\to$ 音符量化 (Quantize 1/16)...」把选中的音符吸附到 1/16 网格；
「编辑 $\to$ 音符人性化 (Humanize)...」打开一个对话框，可设时间抖动、力度
抖动、强度与种子等参数，让机械的量化结果更自然。两者都只作用于当前选区。

\screenshot{ch04-humanize-dialog}{人性化对话框，显示时间抖动、力度抖动与强度输入}

\subsection{音阶与和弦工具}

当前版本以音高车道的键标签提供音高参照；独立的音阶/和弦生成工具是否提供，
以你所用版本的实际界面为准，未出现的入口不要当作已提供。「编辑 $\to$ 鼓机生成
(Drums)...」则确实存在，可按风格、复杂度、种子生成鼓片段。

\section{步进编辑与和弦辅助}

「步进编辑」在这里指按固定步长逐格写音符的做法：用绘制工具在目标格上点出
音符，再靠量化把时间收齐。它不是一键生成整段旋律的功能。

对于和弦，当前可做的是在卷帘里把同一时刻的多个音高画成分层音符，再用试听
核对；自动和弦进行生成不在本章描述范围，未提供的能力不以「已有」表述。

\section{帮助主题映射}

\begin{table}[htbp]
\centering
\caption{本章小节 $\to$ 应用内帮助锚点}
\label{tab:ch04-help}
\begin{tabular}{ll}
\hline
本章小节 & 帮助主题 \\
\hline
打开卷帘 / 两个维度 & \texttt{roll-open-view} \\
缩放与取景 & \texttt{roll-zoom-framing} \\
画音符与选择 & \texttt{roll-note-entry} \\
选择与拖动（单音符边界） & \texttt{roll-note-drag} \\
力度泳道 & \texttt{roll-velocity-lane} \\
量化与人性化 & \texttt{roll-quantize-humanize} \\
音阶与和弦工具 / 步进编辑 & \texttt{roll-scale-step} \\
\hline
\end{tabular}
\end{table}

\paragraph{本章学完能做什么} 在卷帘里画一段旋律、改力度、量化到 1/16，并说清当前拖动只作用于单颗音符。

\paragraph{动手练习} 在一条 MIDI 轨的片段上打开卷帘，用数字键 2 画一小节 C 大调音阶，在力度泳道把第一个音改强，用「音符量化」收到 1/16，再按空格试听；最后试着拖其中一颗音符，确认其余音符不动。
